\subsection{Normalizers}

\begin{lema}\label{lem:normalizadores} Let $\rho=1-x_*=q/n$. The following hold uniformly in $1\le h\le K$. \begin{enumerate}[(i)]
\item $W_R(0)=n^2\rho/\lambda_\rho\,(1+O(n^{-1}))\asymp x_*^2n^2$ and $W_V(0)=H^4(1+O(n^{-1}))$. \item $\sum_v\log
u_v=n\sum_f\int_{I_f}\log u_f+O(1)$. \item $\E_{\PP_h}\sum_t\log W_{\col(t)}(t)\le q\log W_R(0)+H\log W_V(0)-4N+O(n^{2/3})$.
\end{enumerate} \end{lema}

\begin{proof} (i) By quadrature for the rescaled factors, which are positive and Lipschitz, $W_R(0)=x_*^2n^2\int
r_Rr_Cu_Du_A\,(1+O(n^{-1}))=x_*^2n^2\rho/b\,(1+O(n^{-1}))$; moreover $x_*^2/b=1/\lambda_\rho$. In each family the sum of the weights is
$H+O(1)$. For the rows we sum $n-1$ dummy contributions and $\sum u_R=nx_*+O(x_*)$. For the diagonals the sum is $\sum u_D=n\int
u_D+O(1)=n(2-\rho)+O(1)$. Since $W_V(0)$ is the product of the four sums, $W_V(0)=H^4(1+O(n^{-1}))$.

(ii) For the rows, $\sum\log u_R=n\log x_*+\sum\log r_R(\bar r)$. The constant $\log x_*$ is summed exactly $n$ times, and $\log r_R$
is uniformly Lipschitz, so the quadrature error is $O(1)$ independently of $1/x_*$. The same holds for the columns, diagonals, and
antidiagonals. The dummy vertices contribute zero.

(iii) Fix a completed matching and a piece $p$ revealed at time $\mathsf t_p\in[0,1]$. Let $z=1-\mathsf t_p$, and let $W^{(p)}$ denote
the normalizer for the color of $p$ at the time it is revealed. Conditional on the matching and on $\mathsf t_p$, each of the other
pieces independently remains unrevealed with probability $z$. An \emph{ordinary} edge of the color of $p$ has its four vertices covered
by four distinct pieces other than $p$ and is available with probability $z^4$. By the codegree-one property there are $O(n)$
exceptional real edges of weight $O(x_*^2)$. Except for the one contained in $p$, each requires at least one other piece to remain
unrevealed and is available with probability at most $z$. There are $O(n^3)$ exceptional virtual edges, each of weight at most one.
Thus, uniformly in the matching, \[
 \E W^{(p)}\le W_R(0)z^4+Cx_*^2nz+Cx_*^2\ \ (p\text{ real}),\qquad
 \E W^{(p)}\le W_V(0)z^4+Cn^3z+C\ \ (p\text{ virtual}).
\] By Jensen's inequality, (i), and $\log(1+u+v)\le\log(1+u)+\log(1+v)$, the error terms are of orders $1/(nz^3)$ and $1/(n^2z^4)$ for
real $p$ and $1/(nz^3)$ and $1/(n^4z^4)$ for virtual $p$. Since $\int_0^1\log z^4\,dz=-4$ and, for fixed $c>0$,
$\int_0^1\log(1+cn^{-1}z^{-3})\,dz=O(n^{-1/3})$, $\int_0^1\log(1+cn^{-2}z^{-4})\,dz=O(n^{-1/2})$, and
$\int_0^1\log(1+cn^{-4}z^{-4})\,dz=O(n^{-1})$, we obtain \[
 \E\log W^{(p)}\le\log W_{\col(p)}(0)-4+O(n^{-1/3}).
\] Since $\sum_t\log W_{\col(t)}(t)=\sum_p\log W^{(p)}$, summing over the $N$ pieces gives (iii). \end{proof}

\begin{lema}\label{lem:retirada}
$N_h(n)\ge\binom nh Q(n)/h!$.
\end{lema}

\begin{proof} Removing $h$ queens from a solution gives a placement in $\Om_h$, and each placement in $\Om_h$ has at most $h!$
completions, because the $h$ missing queens occupy a permutation of the $h$ empty crossings. Hence $\binom nhQ\le h!N_h$. \end{proof}

\begin{teorema}[Entropy comparison]\label{thm:comparacion} Fix $K$. There exists $d_n=\Delta_n+C_Kn^{2/3}+O_K(\log n)=o(n)$ such that
for all $0\le h\le K$, \[
 \KL{\PP_h}{\QQ_h}\le d_n,\qquad \QQ_h(S)\ge e^{-d_n},\qquad \KL{\mu_h}{\nu'_h}\le d_n .
\] \end{teorema}

\begin{proof}[Proof for $1\le h\le K$] By Lemmas~\ref{lem:identidad} and~\ref{lem:normalizadores} and Stirling's formula ($\log
H!=H\log H-H+O(\log n)$, and likewise for $q!$), \[
 \KL{\PP_h}{\QQ_h}\le q\log\frac{n^2\rho}{\lambda_\rho}-4q-n\sum_f\int\log u_f-(q\log q-q)
 -\log N_h+Cn^{2/3}+O(\log n),
\] where we have used $4H\log H-4\log H!=4H+O(\log n)$ and $N-H=q$. With $q=n\rho$, we have $q\log(n^2\rho)-q\log q=q\log n$, so the
right-hand side is $q\log n+n[-3\rho-\rho\log\lambda_\rho-\sum_f\int\log u_f]-\log N_h+Cn^{2/3}+O(\log n)$. The expression in brackets
is $s^*(\rho)$, since the Euler equation gives $\int\nu\log\nu=\rho\log\lambda_\rho+\sum_f\int m_f\log u_f$ and $\int m_f\log
u_f=\int\log u_f-\int u_f\log u_f$. Thus \[
 \KL{\PP_h}{\QQ_h}\le q\log n+ns^*(q/n)-\log N_h+Cn^{2/3}+O(\log n).
\] By Lemma~\ref{lem:borde}, $ns^*(q/n)=ns_1+2h\log n+O_K(1)$. By Lemma~\ref{lem:retirada}, $-\log N_h\le-\log Q-\log\binom nh+\log
h!=-\log Q-h\log n+O_K(1)$. The terms in $\log n$ cancel exactly ($-h+2h-h=0$), and by~\eqref{eq:Delta},
$\KL{\PP_h}{\QQ_h}\le\Delta_n+C_Kn^{2/3}+O_K(\log n)=d_n$. No rate is needed for $\Delta_n$. Since $\PP_h(S)=1$, Lemma~\ref{lem:exito}
gives $\QQ_h(S)\ge e^{-d_n}$ and $\KL{\PP_h}{\QQ_h(\cdot\mid S)}\le d_n$. Projecting onto the final placement and applying data
processing gives $\KL{\mu_h}{\nu'_h}\le d_n$. \end{proof}
